\scanpage{334}\mspage{250}
% l. 1–4
Ainsi, on est réduit à prouver 4.10.3. Appliquant 4.9.6. $2^\circ$, on est ramené à prouver, pour le foncteur envisagé, les conditions $\mathrm F1$ et $\mathrm F3_{\unsure{\pi|\kappa}}$ (page 225).

% l. 5–7
Dans les arguments qui suivent, on va utiliser les conditions a), b) de 4.10.3, mais non, pour l'instant, la condition c.

% l. 8–15
\uline{Lemme} 4.10.5. Sous les conditions de 4.10.3 a) b) (à l'exclusion de c), pour que (4.10.3.1) soit cofinal et $\Gamma(\unsure{\mathcal E|\mathscr E}/J)$ grde devant $\pi$, il faut et il suffit que le foncteur envisagé satisfasse à la condition $\mathrm F1$ \deleted{\unsure{et à sa variante}} \inserted{rappelée ci-dessous,} et à la variante $\mathrm F'1$ de celle-ci \unsure{énoncée|donnée} ci-dessous :

% l. 16–19
$\unsure{\bullet|\cdot}\ \mathrm F1)$ Pour \unsure{tout|tt} système $(X(j))_{j\in\operatorname{Ob}J}$ d'éléments $X(j)\in\operatorname{Ob}\mathcal E_j$, existe une \uline{section} $\unsure{G|E}$ de $\mathcal E$ sur $J$, et un système de flèches
\[
 X(j)\longrightarrow G(j)\qquad (j\in\operatorname{Ob}J). \tag*{(4.10.\unsure{5|3}.1)}
\]

% l. 20–21
$\mathrm F'1)$ Sous les conditions de $\mathrm F1$ \unsure{précédent|précédentes}, supposons de plus une section $F$ de $\mathcal E$ sur $J$,
% archive sheet number: 333
